\section{Applications of dimension theory}
\label{section-applications-dimension-theory}

\noindent
We can use the results on dimension to prove certain rings
have infinite spectra and to produce more Jacobson rings.

\begin{lemma}
\label{lemma-Noetherian-local-domain-dim-2-infinite-opens}
Let $R$ be a Noetherian local domain of dimension $\geq 2$.
A nonempty open subset $U \subset \Spec(R)$ is
infinite.
\end{lemma}

\begin{proof}
To get a contradiction, assume that $U \subset \Spec(R)$ is finite.
In this case $(0) \in U$ and $\{(0)\}$ is an open subset of $U$ (because
the complement of $\{(0)\}$ is the union of the closures of the other points).
Thus we may assume $U = \{(0)\}$.
Let $\mathfrak m \subset R$ be the maximal ideal.
We can find an $x \in \mathfrak m$, $x \not = 0$ such that
$V(x) \cup U = \Spec(R)$. In other words we see that
$D(x) = \{(0)\}$. In particular we see
that $\dim(R/xR) = \dim(R) - 1 \geq 1$, see Lemma \ref{lemma-one-equation}.
Let $\overline{y}_2, \ldots, \overline{y}_{\dim(R)} \in R/xR$ generate
an ideal of definition of $R/xR$, see Proposition \ref{proposition-dimension}.
Choose lifts $y_2, \ldots, y_{\dim(R)} \in R$, so that
$x, y_2, \ldots, y_{\dim(R)}$ generate an ideal of definition in $R$.
This implies that $\dim(R/(y_2)) = \dim(R) - 1$ and
$\dim(R/(y_2, x)) = \dim(R) - 2$, see
Lemma \ref{lemma-elements-generate-ideal-definition}.
Hence there exists a prime
$\mathfrak p$ containing $y_2$ but not $x$. This contradicts
the fact that $D(x) = \{(0)\}$.
\end{proof}

\noindent
The rings $k[[t]]$ where $k$ is a field, or the ring of $p$-adic
numbers are Noetherian rings of dimension $1$ with finitely many primes.
This is the maximum dimension for which this can happen.

\begin{lemma}
\label{lemma-Noetherian-finite-nr-primes}
A Noetherian ring with finitely many primes has dimension $\leq 1$.
\end{lemma}

\begin{proof}
Let $R$ be a Noetherian ring with finitely many primes.
If $R$ is a local domain, then the lemma follows from
Lemma \ref{lemma-Noetherian-local-domain-dim-2-infinite-opens}.
If $R$ is a domain, then $R_\mathfrak m$ has dimension $\leq 1$
for all maximal ideals $\mathfrak m$ by the local case.
Hence $\dim(R) \leq 1$ by Lemma \ref{lemma-dimension-height}.
If $R$ is general, then $\dim(R/\mathfrak q) \leq 1$
for every minimal prime $\mathfrak q$ of $R$.
Since every prime contains a minimal prime
(Lemma \ref{lemma-Zariski-topology}), this implies $\dim(R) \leq 1$.
\end{proof}

\begin{lemma}
\label{lemma-finite-type-algebra-finite-nr-primes}
Let $S$ be a nonzero finite type algebra over a field $k$.
The following are equivalent
\begin{enumerate}
\item $\dim(S) = 0$,
\item $S$ has finitely many primes,
\item $S$ has finitely many maximal ideals,
\item $\Spec(S)$ satisfies one of the equivalent conditions of
Lemma \ref{lemma-ring-with-only-minimal-primes},
\item $\dim_k(S) < \infty$,
\item $S$ is Artinian,
\item $\Spec(S)$ is a discrete topological space,
\item add more here.
\end{enumerate}
\end{lemma}

\begin{proof}
It is immediate from the definitions that (1) is equivalent to (4)
by looking at part (5) of Lemma \ref{lemma-ring-with-only-minimal-primes}.
Recall that $\Spec(S)$ is sober, Noetherian, and Jacobson, see
Lemmas \ref{lemma-spec-spectral}, \ref{lemma-Noetherian-topology},
\ref{lemma-finite-type-field-Jacobson}, and \ref{lemma-jacobson}.
If $S$ has dimension $0$, then every point defines an
irreducible component and there are only a finite number
of irreducible components (Topology, Lemma \ref{topology-lemma-Noetherian}).
Thus (1) implies (2). Trivially (2) implies (3).
If (3) holds, then $\Spec(S)$ is discrete by
Topology, Lemma \ref{topology-lemma-finite-jacobson}
and hence the dimension of $S$ is $0$.

\medskip\noindent
At this point we know that (1) -- (4) are equivalent.
The implication (5) $\Rightarrow$ (6) is
Lemma \ref{lemma-finite-dimensional-algebra}.
The implication (6) $\Rightarrow$ (7) follows from
Proposition \ref{proposition-dimension-zero-ring}.
The implication (7) $\Rightarrow$ (4) is immediate.
Conversely, if $S$ satisfies (1) -- (4), then $S$ has finitely
many primes $\mathfrak m_1, \ldots, \mathfrak m_r$ all maximal.
Note that $\kappa(\mathfrak m_i)$ is a finite extension of $k$ by the
Hilbert Nullstellensatz (Theorem \ref{theorem-nullstellensatz}).
By Proposition \ref{proposition-dimension-zero-ring} we also see
that $S$ is Artinian. Next, Lemma \ref{lemma-artinian-finite-length}
tells us that $\text{length}_S(S) < \infty$. Thus $\dim_k(S) < \infty$ by
Lemma \ref{lemma-pushdown-module}. We conclude that
(1) -- (7) are equivalent. (Note: another and more standard
way to prove $\dim(S) = 0 \Rightarrow \dim_k(S) < \infty$
is to use Noether normalization, but we don't have this available to us yet.)
\end{proof}

\begin{lemma}
\label{lemma-noetherian-dim-1-Jacobson}
Noetherian Jacobson rings.
\begin{enumerate}
\item Any Noetherian domain $R$ of dimension $1$
with infinitely many primes is Jacobson.
\item Any Noetherian ring such that every prime
$\mathfrak p$ is either maximal or contained in
infinitely many prime ideals is Jacobson.
\end{enumerate}
\end{lemma}

\begin{proof}
Part (1) is a reformulation of Lemma \ref{lemma-pid-jacobson}.

\medskip\noindent
Let $R$ be a Noetherian ring such that
every non-maximal prime $\mathfrak p$ is contained
in infinitely many prime ideals.
Assume $\Spec(R)$ is not Jacobson to get
a contradiction.
By Lemmas \ref{lemma-irreducible}
and \ref{lemma-Noetherian-topology}
we see that $\Spec(R)$ is a sober, Noetherian topological space.
By Topology, Lemma \ref{topology-lemma-non-jacobson-Noetherian-characterize}
we see that there exists a non-maximal ideal $\mathfrak p \subset R$
such that $\{\mathfrak p\}$ is a locally closed subset of
$\Spec(R)$. In other words, $\mathfrak p$ is not maximal
and $\{\mathfrak p\}$ is an open subset of $V(\mathfrak p)$.
Consider a prime $\mathfrak q \subset R$ with
$\mathfrak p \subset \mathfrak q$. Recall that the topology on the spectrum of
$(R/\mathfrak p)_{\mathfrak q} = R_{\mathfrak q}/\mathfrak pR_{\mathfrak q}$
is induced from that of $\Spec(R)$, see Lemmas
\ref{lemma-spec-localization} and \ref{lemma-spec-closed}.
Hence we see that $\{(0)\}$ is a locally closed subset of
$\Spec((R/\mathfrak p)_{\mathfrak q})$. By
Lemma \ref{lemma-Noetherian-local-domain-dim-2-infinite-opens}
we conclude that $\dim((R/\mathfrak p)_{\mathfrak q}) = 1$.
Since this holds for every $\mathfrak q \supset \mathfrak p$
we conclude that $\dim(R/\mathfrak p) = 1$. At this point we use
the assumption that $\mathfrak p$ is contained in infinitely many
primes to see that $\Spec(R/\mathfrak p)$ is infinite.
Hence by part (1) of the lemma we see that
$V(\mathfrak p) \cong \Spec(R/\mathfrak p)$
is the closure of its closed points.
This is the desired contradiction since it means that
$\{\mathfrak p\} \subset V(\mathfrak p)$ cannot be open.
\end{proof}




















